\documentclass[11pt]{article}

\usepackage[margin=1in]{geometry}
\usepackage[T1]{fontenc}
\usepackage{lmodern}
\usepackage{amsmath,amssymb,amsthm}
\usepackage{booktabs}
\usepackage{microtype}
\usepackage{hyperref}

\newtheorem{theorem}{Theorem}
\newtheorem{remark}{Remark}

\title{Conservation is the only exact content: a verification layer for climate models}
\author{Rob Mitchell}
\date{}

\begin{document}
\maketitle

\begin{abstract}
Climate projections are the most consequential predictions we make, and the
field has no verification layer that distinguishes a \emph{verified} projection
from an \emph{extrapolated} one. We apply the PACLAND verification discipline ---
the two-level certificate, the verification deficit, and the constraint-first
decomposition --- to the CMIP6 ensemble. Three results. \textbf{(1) The leak:}
12 of 19 models fail a three-ledger conservation guard (energy, mass, salt), and
the leak is localized to the atmosphere--ocean coupling, not the ocean.
\textbf{(2) The fix:} drift-correction --- the field's own post-processing repair
--- restores conservation in every model (12 $\to$ 0 failures), but a
retrospective test against HadCRUT5 shows it does \emph{not} fix the temperature
bias; that bias comes from sensitivity, and the sensitivity+forcing gate fixes it
(ensemble ECS 3.78 $\to$ 3.08~K, matching the observed 3.0~K). \textbf{(3) The
theorem:} conservation is the \emph{only} exact, model-independent content of a
climate model --- the climate analog of the flux tadpole --- and every predictive
quantity (ECS, forcing, the 2100 warming) is a model-dependent range, not an exact
observable. The verification deficit climbs to $m = 7$ (4 empirical observables +
3 exact conservation ledgers) and is \emph{saturated}: the constraint contribution
is capped at 3, exactly as the flux sector is capped at the tadpole.
\end{abstract}

\section{The verification deficit, applied to climate}

The deficit counts observables that are simultaneously \emph{measured, computable,
and validated}. On the historical record, $m = 4$ (GMST, EEI, OHC, GMSL --- a
$\approx 68\%$ certificate radius). The 2100 projection has $m = 0$: no future
state is measured, so the projection is extrapolation, not verification.

\section{The leak --- three conservation ledgers}

A conserving model must close three independent ledgers in its control run:
\emph{energy} (net top-of-atmosphere flux $\approx 0$), \emph{mass} (ocean mass
$\approx 0$), \emph{salt} (salinity drift $\approx 0$). Measured (Irving et
al.~2021, 19 models): 12 leak energy, 6 leak mass, 2 leak salt. Salt --- the
purely-ocean ledger --- is conserved orders of magnitude better than energy and
mass, which cross the air--sea boundary; the leak therefore lives in the
\emph{atmosphere--ocean coupling}, not the ocean model. The worst model
(CNRM-ESM2-1) leaks at $\approx 1.9\times$ the Earth's energy imbalance.

Widening to all 53 models (a raw second-source recomputation of the piControl
net-TOA, validated against Irving within $0.03$~W/m$^2$ on the overlapping models)
finds \emph{22 of 27} previously-unmeasured models leak energy --- the models
Irving never covered are, if anything, leakier.

\section{The fix --- drift-correction}

The field's own repair is drift-correction: subtract the measured leak. Applied to
all models, every one conserves by construction (12 $\to$ 0 failures). But a
retrospective test against HadCRUT5 (HadGEM3-GC31-LL, historical + piControl)
shows the correction does \emph{not} improve the temperature match (RMSE
$0.27 \to 0.29$~K; end warming $+1.92$ vs $+0.93$~K observed). The reason is
stated, not hidden: HadGEM3's $\sim 1$~K over-warming comes from its
\emph{sensitivity} (a ``hot'' model), not its conservation drift
($0.02$~K/century). Two tools, two jobs: drift-correction fixes
\emph{conservation}; the sensitivity+forcing gate fixes the \emph{temperature
bias} (ensemble ECS $3.78 \to 3.08$~K, matching the observed $3.0$~K). Three
independent observational sources agree: Sherwood et al.~2020 (ECS), HadCRUT5
(temperature), Loeb et al.~2021 (EEI).

\section{The Climate Constraint-First Theorem}

\begin{theorem}
The only exact, model-independent content of a climate model is its conservation
laws (energy, mass, salt). Every predictive quantity --- ECS, forcing, the 2100
warming --- is a model-dependent range, not an exact observable.
\end{theorem}

\textbf{Arm A (exact).} Conservation is an identity, not an approximation: a model
that does not close the ledgers is broken, not ``approximately wrong''.
\textbf{Arm B (non-exact).} The 53 CMIP6 models span ECS $1.83$--$5.64$~K (a
range, not a value); any two disagreeing models falsify ``ECS is exact''; the
projection has $m = 0$. This is the climate instance of the one spine:
\emph{``don't ask a complex system for its values --- ask for its constraints.''}

\section{The deficit, climbed and capped}

$m = 7$ --- 4 empirical observables + 3 exact conservation ledgers --- with the
ceiling proved by the theorem: conservation is the only exact content, so the
constraint contribution is \emph{saturated at 3}, exactly as the flux sector is
saturated at 1 (the tadpole). Further deficit growth is empirical only.

\section*{Reproduce}

\noindent
\texttt{PYTHONPATH=src python3 scripts/validate\_climate\_constraint\_first.py}\\
\texttt{PYTHONPATH=src python3 -m pacland.climate.drift\_correction}\\
\texttt{PYTHONPATH=src python3 scripts/climate\_retrospective.py}

\noindent Code: \texttt{src/pacland/climate/} \quad data: \texttt{data/climate/}
\quad 305 tests green.

\end{document}
